\documentclass{article}

\usepackage[preprint]{colm2026_conference}

\usepackage{amsmath}
\usepackage{microtype}
\usepackage{hyperref}
\usepackage{url}
\usepackage{booktabs}
\usepackage{graphicx}
\usepackage{lineno}

\definecolor{darkblue}{rgb}{0, 0, 0.5}
\hypersetup{colorlinks=true, citecolor=darkblue, linkcolor=darkblue, urlcolor=darkblue}

\title{When Self-Consistency Backfires: Majority Vote Hurts the Majority of Hard Science Problems for Small LLMs}

\author{Utkarsh Bahuguna \\
Scaler School of Technology \\
\texttt{utkarshbahuguna10@gmail.com}}

\newcommand{\fix}{\marginpar{FIX}}
\newcommand{\new}{\marginpar{NEW}}

\begin{document}

\ifcolmsubmission
\linenumbers
\fi

\maketitle

\begin{abstract}
Self-consistency (SC) via majority vote reduces per-problem accuracy on most
GPQA Diamond problems for small instruction-tuned models: 56.6\% of problems
for Qwen2.5-7B and 65.7\% for Llama-3-8B. The obvious remedy is a
verifier-free confidence gate that decides which problems to vote on. This
version reports that the most natural repair to that remedy also fails, and
separates three signal failures that v1 treated as one. A token-entropy gate
fails for a measurement reason rather than a substantive one: a per-token
average over a chain averaging 602 tokens is dominated by fluency, and
measured, the chain average sits 0.0005 nats from the average over non-answer
tokens and 0.1695 nats from the answer token's own value. We therefore measure
confidence at the answer span itself, where that dilution cannot apply. On
Qwen2.5-7B-Instruct-Turbo, 198 problems at 64 samples each, a sample whose
answer contradicts its own problem's plurality still emits that answer at a
median margin of 20.52 nats, with 75.7\% of such samples above 10 nats. Both
quantities were pre-registered with thresholds fixed in advance and tested once
on 69 problems that no exploratory analysis had read; both passed. The unit is
the whole result: pooled across the benchmark the same margin does separate
correct from incorrect samples, by $+0.0604$ on the fraction above 10 nats
$[+0.0183, +0.1017]$, excluding zero; measured within each problem and paired
across its own samples it does not, at $-0.0168$ $[-0.0527, +0.0182]$, crossing
zero. So the claim is not that token log-probabilities carry no information,
but that a signal with genuine across-question discrimination is close to
useless for the within-question decision self-consistency actually poses. The
plurality-agreement gate's failure, unlike the other two, remains without a
mechanism, and we report it as an open problem. These new claims rest on one
model: a registered second-model replication was sampled and could not be
evaluated, and we report that rejection rather than the result. We separately
report that on hosted serverless inference at a small budget, three
reasoning-native models could not be evaluated at all, for three distinct and
separately measured reasons; all three are downloadable, so this bounds what a
metered per-token API buys rather than what is knowable.
\end{abstract}

\section*{Changes from v1}

\textbf{Peer-review status.} \textbf{The COLM 2026 Workshop on Efficient
Reasoning accepted v1 of this paper} (arXiv:2608.11403), whose content is
the original Introduction, Setup, Results and
Discussion. \textbf{No reviewer saw the material added in v2}: the
answer-token margin and the commitment result, the reasoning wall, the
rejected second-model replication, the reconstructed estimator, and the
corrections section. v2 also \emph{corrects} the accepted version:
Section~\ref{sec:why} replaces the accepted paper's mechanism section, and one
Discussion sentence is withdrawn as stated. Readers relying on peer review
should treat the accepted claims and the new ones differently, and the list
below marks which is which.

For readers who cited the original. \textbf{The headline results are
unchanged: no v1 result was recomputed, re-estimated, or withdrawn.} What
moved is the explanation of why the gates fail.

\begin{itemize}
\item \textbf{Corrected.} Section~\ref{sec:why}'s v1 thesis, ``confidence does
not track correctness'', inferred one mechanism from two gates failing
together. The failures have different causes and one of them is a measurement
artifact.
\item \textbf{Corrected, narrowed to a unit.} v1's Discussion claim that
``confidence on these problems does not indicate correctness'' is false as
stated across questions: pooled, the answer-token margin separates correct
from incorrect by $+0.0604$ $[+0.0183, +0.1017]$. Within a problem, paired, it
does not: $-0.0168$ $[-0.0527, +0.0182]$. The defensible claim is the
within-question one.
\item \textbf{Partly answered, and rescoped.} v1's ``we do not test
reasoning-native models, which we flag as the central open question'' is no
longer open in the same way: on hosted serverless inference at a small budget
it is unmeasurable (Section~\ref{sec:wall}). It remains open for the
open-weights route.
\item \textbf{Separated.} The entropy gate's failure was grouped with
agreement's under one known-mechanism heading. It was reading a diluted
statistic, which is miscounted evidence rather than evidence about confidence.
\item \textbf{Added.} The answer-token margin (the ``richer signal'' v1's
Limitations named as future work), the reasoning wall, a registered
replication that this paper's own comparability rule rejected, a reconstructed
few-sample estimator, and a corrections section.
\end{itemize}

\textbf{v2 closes a limitation v1 stated about itself.} v1 wrote that
evaluating final-answer margins ``requires per-token log-probability arrays,
which our pipeline did not store''. This version stores them and evaluates that
signal. It does not behave differently in the way that sentence hoped.

\section*{Data availability}

Derived, question-free data reproducing every number in this paper is archived
open access at \href{https://doi.org/10.5281/zenodo.21933418}{10.5281/zenodo.21933418} (CC BY 4.0). The raw per-sample
store, including model chains of thought, is archived at
\href{https://doi.org/10.5281/zenodo.21933422}{10.5281/zenodo.21933422} as a \textbf{restricted, request-access} record
on the same terms as GPQA~\citep{rein2023gpqa}; it is not open data, for the
reason given in Section~\ref{sec:cot-release}.

\section{Introduction}

Inference-time compute is a primary lever for LLM reasoning, and
self-consistency---sample several chains of thought, return the plurality
answer---is among the simplest and most widely used techniques in this
family~\citep{wang2023self}. It is commonly assumed to be a low-risk
accuracy boost: sample more, vote, do at least as well.

That assumption fails on hard problems. When a model places its highest
probability on an incorrect answer across independent samples, more samples
only entrench the wrong vote. We call this \emph{backfire}:
$\text{mv\_gain} = \text{MV\_acc}(64) - \text{MV\_acc}(1) < 0$, where
$\text{MV\_acc}(N)$ is the expected accuracy of a majority vote over $N$
samples. We then ask
the question a cost-conscious practitioner would ask: can a cheap,
verifier-free signal computed from a few samples tell you which problems to
vote on and which to skip? We test the two most natural signals, plurality
agreement and token-level entropy, and find both fail, and we identify why.

That self-consistency can hurt is not itself new, and the underlying
miscalibration is well documented~\citep{guo2017calibration,
kadavath2022language}. Our contribution is to quantify, pre-register, and
confirm how often it happens on a full hard benchmark, across two model
families, and to show that two cheap gates cannot prevent it.

\textbf{Contributions:}
\begin{itemize}
  \item On all 198 GPQA Diamond problems, majority voting backfires on the
        majority of problems for two model families (56.6\% and 65.7\%), with
        Qwen the primary demonstration and Llama corroborating from a
        near-chance baseline. The effect was pre-registered on a 151-problem
        held-out split and confirmed (all four hypotheses passed).
  \item We show the per-problem routing headroom is real: a grid oracle (best
        $N$ per problem across $\{1,2,4,8,16,32,64\}$) marks a 14 to 17 point
        ceiling above $N=1$. Neither a plurality-agreement gate nor a
        token-entropy gate reaches it, and neither moves accuracy more than
        0.002 from fixed-budget voting at $N=64$.
  \item We give the mechanism: agreement does not track correctness, and for
        the weaker model higher agreement does not indicate higher accuracy.
\end{itemize}


\textbf{New in v2.} v1 tested two such signals, plurality agreement and
token-level entropy, found both fail, and offered one explanation for both.
This version tests a third, the answer-token log-probability margin, which
v1's own Limitations named as the obvious next candidate and could not
evaluate because its pipeline stored only a mean entropy scalar. We collect
the per-token arrays, measure the margin at the answer span, and report that it
also fails. The three failures do not share an explanation, and separating them
is the main interpretive change in this version: one is a measurement artifact,
one has a mechanism we can state and test, and one remains unexplained.

\section{Setup}

\textbf{Dataset and design.} We use the full GPQA Diamond
benchmark~\citep{rein2023gpqa}, 198 graduate-level multiple-choice questions
in biology, chemistry, and physics. We adopt a pre-registered confirmatory
design: 47 problems are \textbf{exploratory} (the hypotheses below were
generated from them), and the remaining 151 are a \textbf{confirmatory} test
set. The hypotheses and thresholds were locked and git-tagged before any
confirmatory analysis, at tag \texttt{backfire-prereg-v1.0} in the repository
at \url{https://github.com/u7k4rs6/self-consistency-backfire}. We report
exploratory, confirmatory, and pooled (198) values throughout, but decide
PASS/FAIL on the confirmatory set only.

\textbf{Models.} Two instruction-tuned models from different families, served
on Together AI: Qwen2.5-7B-Instruct-Turbo and Meta-Llama-3-8B-Instruct-Lite.
Both are small (7 to 8B) and non-reasoning. $N=64$ samples per problem,
temperature 0.7, single locked prompt template, byte-identical across all runs
as verified by SHA-256. Five-pass answer extraction; parse rate 99.5\% (Qwen)
and 98.6\% (Llama). Llama has exactly 64 stored samples per problem. For Qwen
the 151 confirmatory problems have exactly 64 at this temperature and the 47
exploratory problems have 65 to 72, retained from earlier sampling runs, so
pool size and split membership coincide exactly; three further problems hold
side-test samples at other temperatures, which the analysis filters out.
$\text{MV\_acc}(N)$ draws subsets of size $\min(N, n)$ without replacement from
all $n$ stored samples for a problem, so for those 47 the estimate at $N=64$ is
taken over a larger pool than for the rest. Reclassifying every problem from
its first 64 samples leaves the backfire count unchanged at 112 of 198, with no
problem crossing the backfire boundary
(\texttt{scripts/verify\_truncation\_invariance.py}).

\textbf{Metrics.} $\text{MV\_acc}(N)$ is expected majority-vote accuracy over
$N$ samples, ties broken uniformly at random independent of ground truth,%
\footnote{This v1 sentence is retained verbatim. The implementation breaks ties
lexicographically rather than at random, and the published numbers follow the
implementation; see disclosure 9 in Section~\ref{sec:corrections}, which
quantifies the difference as a third-decimal effect that changes no
conclusion.}
estimated by Monte Carlo over subsets. Backfire is $\text{mv\_gain} < 0$.

\textbf{Routing and gates.} The \textbf{oracle} routes each problem to the
$N$ that maximizes majority-vote accuracy across $\{1,2,4,8,16,32,64\}$ using
ground truth; it is a theoretical upper bound, not a deployable method. The
\textbf{agreement gate} returns the probe plurality when its fraction over $k$
samples is at least $\tau$, else votes at $N=64$. The \textbf{entropy gate}
returns the probe plurality when the mean per-token entropy over $k=4$ probe
samples falls below a threshold, else votes at $N=64$. Both gates are
verifier-free; ground truth is used only to score the result. Logprobs are available for all
151 confirmatory problems but not for the 47 older exploratory ones, so the
entropy gate is evaluated on the confirmatory set. Its threshold is selected
on that same set by an argmax over 40 candidate values, maximizing capture
against the binary $\{N=1, N=64\}$ oracle relative to an
$\text{MV\_acc}(64)$ baseline (Qwen: 0.113, Llama: 0.139).

\textbf{Uncertainty.} 95\% confidence intervals are problem-level bootstrap
(1000 iterations, seed 42).


\subsection{New in v2: the answer-token margin}
\label{sec:margin-setup}

v1 retained only a mean per-token entropy scalar, which is why it could not
evaluate a margin. This version re-samples Qwen2.5-7B-Instruct-Turbo on all 198
problems at $M=64$, cap 2048, requesting log-probabilities at depth 5, and
stores the full response object: 12,672 samples.

The \textbf{answer-token margin} is the log-probability of the emitted option
letter minus that of the highest-scoring alternative option letter, at the
answer position. The provider returns the top $k$ alternatives per token,
$k=5$ in practice, and five slots need not contain every option letter. With
two or more option letters present the margin is \emph{measured}: the
second-highest present option is a returned value and any absent letter lies at
or below the smallest returned value, so no absent letter can outrank it. With
fewer than two present it is \emph{right-censored} at the top letter minus the
smallest returned log-probability, and recorded as a bound. It is never
imputed: filling a missing letter at the censoring bound would understate the
margin exactly on the problems where the model is most certain, which are the
ones a confidence gate cares about most. Of 12,672 samples, 12,344 margins are
measured, 265 right-censored, and 63 have no margin at all because no answer
was extracted; those 63 are the same samples behind the 0.9950 answer rate and
are never scored as incorrect.

The new run's answer rate is 0.9950 $[0.9936, 0.9961]$ against v1's 0.9946
$[0.9932, 0.9958]$ for the same model, so the two runs describe the same output
population, not merely the same prompts. Exposure was tracked per analysis: 129
problems were read by some exploratory analysis and 69 were not. The claims in
Section~\ref{sec:commitment} were registered with thresholds fixed before those
69 were examined, and were examined once.

\section{Pre-Registered Confirmatory Results}

All four pre-registered hypotheses pass on the 151 confirmatory problems
(Table~\ref{tab:confirmatory}). Thresholds were fixed before any confirmatory
analysis and set with margin on the permissive side of the exploratory point
estimates rather than at them, so each is a genuine prediction rather than a
restatement. PH4 had no exploratory estimate, since logprobs were not collected
for the exploratory split, and was given the same 10\% ceiling as PH2.
PH2 and PH4 were pre-registered against a binary $\{N=1, N=64\}$ oracle; the
grid oracle results in Sections~\ref{sec:oracle} and~\ref{sec:gates} use a
more generous ceiling and are reported separately as post-hoc analysis.

\begin{table}[t]
\begin{center}
\begin{tabular}{llllc}
\toprule
\textbf{Hyp} & \textbf{Prediction (both models)} & \textbf{Qwen2.5-7B} & \textbf{Llama-3-8B} & \textbf{Result} \\
\midrule
PH1 & backfire rate $\geq$ 33\%          & 60.3\% [53.0, 68.2] & 65.6\% [58.3, 73.5] & PASS \\
PH2 & agree gate capture $\leq$ 10\%   & 0.8\% [$-$89.1, 68.1] & $-$1.6\% [$-$92.9, 74.2] & PASS \\
PH3 & top-agree-bin acc.\ $\leq$ 70\%  & 51.2\% ($n$=43)$^{\dagger}$ & 14.3\% ($n$=21)$^{\dagger}$ & PASS \\
PH4 & entropy gate capture $\leq$ 10\% & 0.5\%$^{\dagger}$    & 0.9\%$^{\dagger}$      & PASS \\
\bottomrule
\end{tabular}
\end{center}
\caption{Confirmatory hypotheses ($n=151$). Thresholds locked and git-tagged
before analysis. Captures in this table are measured against the binary
$\{N=1, N=64\}$ oracle relative to an $\text{MV\_acc}(64)$ baseline, which is
why they are much smaller than the grid-oracle captures reported in
Section~\ref{sec:gates}. Bracketed values are 95\% bootstrap intervals
(problem-level, 1000 iterations, seed 42); $\dagger$ marks quantities for which
no interval was computed. PH2's intervals are wide enough to be consistent with
capture well above the 10\% threshold: the pre-registered decision rule is on
the point estimate, and we report the intervals alongside it.}
\label{tab:confirmatory}
\end{table}

\textbf{On PH2, added in v2 and leaving the published verdict untouched.} PH2
is recorded above as PASS, which is the correct application of its
pre-registered decision rule: that rule is stated on the point estimate, and
the point estimates are 0.8\% and $-1.6\%$, both below the 10\% ceiling. The
intervals around them are wide enough to be consistent with capture well above
that ceiling. The verdict as published stands and is not revised here; what v2
adds is that the interval, not the point estimate, is the honest measure of
what this hypothesis established, and it established less than the PASS alone
suggests.

\section{Results (Pooled, 198 Problems)}

\subsection{Backfire affects the majority of problems, precisely estimated}

Majority voting reduces per-problem accuracy on most problems for both models
(Figure~\ref{fig:backfire}): the pooled backfire rate is 56.6\% (95\% CI
[49.5, 63.6]) for Qwen and 65.7\% ([59.1, 71.7]) for Llama. Expanding from
47 to 198 problems roughly halved the interval width, and both rates sit well
above the 33\% threshold.

The aggregate and per-problem pictures diverge. Voting barely moves aggregate
accuracy (Qwen 0.342 to 0.369, Llama 0.273 to 0.313) while harming the
majority of problems individually. An average that looks flat or mildly
positive can hide widespread per-problem harm. The worst single problem loses
47 points (Qwen) or 46 (Llama); a few gain as much as 66 or 70. The
asymmetry is real but rare: large gains exist, yet most problems backfire.

\begin{figure}[t]
\begin{center}
\includegraphics[width=\linewidth]{backfire_both.png}
\end{center}
\caption{$\text{mv\_gain} = \text{MV\_acc}(64) - \text{MV\_acc}(1)$ per problem.
Dark bars are backfire ($\text{mv\_gain} < 0$).}
\label{fig:backfire}
\end{figure}

Backfire is not uniform across domains (Table~\ref{tab:domain}). Chemistry is the
dominant contributor: it accounts for 93 of the 198 problems and shows the highest
backfire rate for both models (66.7\% Qwen, 68.8\% Llama). Physics is intermediate
(50.0\% Qwen, 67.4\% Llama), and biology is least affected, though its small problem
count ($n=19$) makes that estimate noisy. We report this breakdown descriptively: the
per-domain counts, biology especially, are too small to support claims about why
backfire concentrates where it does.

\begin{table}[t]
\begin{center}
\begin{tabular}{lcccc}
\toprule
\textbf{Domain} & \textbf{Qwen $n$} & \textbf{Qwen backfire} & \textbf{Llama $n$} & \textbf{Llama backfire} \\
\midrule
Biology   & 19  & 36.8\% & 19  & 42.1\% \\
Chemistry & 93  & 66.7\% & 93  & 68.8\% \\
Physics   & 86  & 50.0\% & 86  & 67.4\% \\
\midrule
Pooled    & 198 & 56.6\% & 198 & 65.7\% \\
\bottomrule
\end{tabular}
\end{center}
\caption{Per-domain backfire rate (pooled 198). Chemistry shows the highest backfire
rate in both models and accounts for nearly half of all problems.}
\label{tab:domain}
\end{table}

\subsection{The oracle upper bound is real but not reachable}
\label{sec:oracle}

A grid oracle that routes each problem to the best majority-vote accuracy
across all $N$ in $\{1,2,4,8,16,32,64\}$ reaches 0.482 (Qwen) and 0.439
(Llama), a ceiling 14 points above $N=1$ for Qwen and 17 points above for
Llama. This is a more generous bound than a binary $\{N=1, N=64\}$ oracle,
because it allows any compute level per problem. It marks how much accuracy a
perfect per-problem routing decision could recover. It is an upper bound, not
a method (Figure~\ref{fig:pareto}). The question this raises is how much of
that ceiling is reachable without ground truth, which
Section~\ref{sec:gates} answers by testing two verifier-free gates.

\begin{figure}[t]
\begin{center}
\includegraphics[width=\linewidth]{pareto_both.png}
\end{center}
\caption{Accuracy vs mean compute: fixed-budget voting, the oracle upper
bound, and the verifier-free agreement and entropy gates. The plotted oracle is
the binary $\{N=1, N=64\}$ oracle (0.462 Qwen, 0.410 Llama), not the more
generous grid oracle whose higher values are quoted in
Section~\ref{sec:oracle}. The entropy series is computed on the confirmatory
151 only, a harder subset than the pooled 198 used for the other series, so its
lower position does not by itself indicate a worse method.}
\label{fig:pareto}
\end{figure}

\subsection{Two verifier-free gates fail to capture it}
\label{sec:gates}

Neither cheap gate recovers the headroom. The agreement gate ($k=8$,
$\tau=0.75$) reaches 0.368 accuracy for Qwen and 0.312 for Llama, against
0.369 and 0.313 for fixed-budget voting at $N=64$, differences of 0.0006 and
0.0014. We report these magnitudes rather than a significance claim, having
computed no paired test. Against that ceiling, the agreement gate captures
18.7\%
of available headroom for Qwen and 23.4\% for Llama (headroom is measured
against the grid oracle of Section~\ref{sec:oracle}, not the binary oracle in
Table~\ref{tab:confirmatory}): non-trivial fractions of a generous ceiling,
yet accuracy remains flat. The gate runs at a mean of 40.4 samples per problem
for Qwen and 45.6 for Llama, which places it between two fixed-budget grid
points. Against $N=64$ it reaches 0.3680 and 0.3117, within 0.0006 and 0.0014
of that baseline, on 63\% and 71\% of the samples. Against $N=32$, which scores
0.3621 and 0.3030, it is higher by 0.0060 and 0.0087 but spends 26\% and 43\%
more samples. What the gate buys relative to a flat $N=64$ is therefore compute
rather than accuracy. Those capture figures are measured against
$\text{MV\_acc}(1)$, so they credit the gate with headroom that voting alone
already recovers. Measured instead against $\text{MV\_acc}(64)$, which isolates
what the routing decision adds over simply voting, the capture is PH2 in
Table~\ref{tab:confirmatory}: 0.8\% for Qwen and $-$1.6\% for Llama.
Agreement routing therefore contributes essentially nothing to accuracy that
fixed-budget voting does not already provide. The gate spends $k$ probe
samples on every problem, then either returns the probe plurality (if the
plurality fraction $\geq \tau$) or proceeds to $N=64$. Every gate-versus-baseline
comparison we report uses fixed-budget voting at a flat $N=64$; we do not
compute a compute-matched baseline.

The entropy gate is evaluated entirely within the confirmatory 151, since
logprobs exist only for that split. Measured
consistently within it (the threshold was also selected on this set, so these
are in-sample figures), the gate captures 0.2\% of the grid-oracle headroom
for Qwen and 31.2\% for Llama, with gate accuracy of 0.318 (Qwen, threshold
0.113) and 0.306 (Llama, threshold 0.139). Because this capture is computed
on a different problem set from the agreement-gate figures above, the two
should not be read as like-for-like: the entropy percentages are
confirmatory-151 values and the agreement percentages are pooled-198 values.

The like-for-like baselines are therefore the confirmatory-151 fixed-budget
accuracies, and the two models differ. For Qwen these are
$\text{MV\_acc}(1) = 0.3182$ and $\text{MV\_acc}(64) = 0.3179$, against a gate
accuracy of 0.3184: the gate matches single-sample accuracy to three decimals,
so it does nothing rather than actively hurting, and the apparent five-point
drop against the 0.369 quoted above is entirely the pooled-versus-confirmatory
mismatch. For Llama they are 0.2518 and 0.3046, against a gate accuracy of
0.3055: there the gate is level with full voting and well above single-sample.
One further point: Qwen's confirmatory $\text{MV\_acc}(64)$ sits marginally
below its $\text{MV\_acc}(1)$, so on this split voting at $N=64$ is very
slightly worse than not voting at all. The gap is small and we do not lean on
it; it complements the pooled aggregate reported above (0.342 to 0.369) by
showing the per-problem backfire effect surfacing faintly in the aggregate.

Mean per-token entropy modestly predicts which problems
backfire (AUC 0.631 for Qwen, 0.523 for Llama), but predictive signal does
not translate into routing accuracy because the gate cannot act on that signal
without misclassifying enough problems to erase the gain. Across all 22
operating points swept ($k \in \{4, 8\}$ and $\tau$ from 0.50 to 1.00 in steps
of 0.05), none beats fixed-budget meaningfully (Figure~\ref{fig:sweep}). The
obvious agreement signal and an uncertainty signal both fail.

\begin{figure}[t]
\begin{center}
\includegraphics[width=\linewidth]{gate_sweep_both.png}
\end{center}
\caption{Agreement-gate accuracy across the $(k, \tau)$ sweep. The lower row
plots the difference from fixed-budget voting at a flat $N=64$. No operating
point wins.}
\label{fig:sweep}
\end{figure}

\subsection{Why: three signals, three different failures}
\label{sec:why}

\textbf{This subsection replaces v1's ``Why: confidence does not track
correctness''.} v1 offered a single account for two failing gates. With a third
signal measured, that account no longer holds as stated: the three failures
have three different statuses, and only one of them is explained by
miscalibration.

\textbf{Signal 1: plurality agreement. Fails, and we cannot say why.} The
evidence is unchanged (Table~\ref{tab:calibration},
Figure~\ref{fig:calibration}). Even the highest-agreement bin is far from
reliable: Qwen's plurality is correct 52.5\% of the time there, and Llama's
only 28.6\%, lower than its own low-agreement bin, and Llama's accuracy is not
monotone in agreement. What v2 changes is the status of the explanation, not
the evidence. Plurality agreement is computed over \emph{answers}, not over
tokens: it never touches a log-probability, so neither dilution nor
answer-token saturation can apply to it, and the account developed below for
the other two signals says nothing about this one. \textbf{Agreement's failure
therefore remains without a mechanism, and we state that as an open problem
rather than absorbing it.} Calling it an instance of general
overconfidence~\citep{guo2017calibration, kadavath2022language} is a label
rather than an account: it does not predict that Llama's top bin should be
\emph{worse} than its bottom one, and it does not say what would have to be
true for agreement to work on some other benchmark. This is the largest
unresolved question in the paper, and it is as unresolved in v2 as in v1. The
difference is that v1's framing made it look answered.

\textbf{Signal 2: mean token entropy. Miscounted evidence, not evidence about
confidence.} The entropy gate averaged a per-token quantity over the whole
chain. In the margin-v1 store that chain averages 602.21 tokens, of which the
answer token is \emph{one}. Measured over 2,560 samples on 40 problems drawn
with a recorded seed, using per-token entropy over the returned top-5
alternatives, the mean over the whole chain is 0.2232 nats, the mean over
non-answer tokens is 0.2237, and the value at the answer token is 0.0536 with a
median of exactly 0.0000. The chain average therefore sits 0.0005 nats from the
non-answer average and 0.1695 nats from the answer token's own value: the
gate's statistic is, to four decimal places, a measurement of the prose. A gate
thresholding it is thresholding fluency. \textbf{This is a measurement artifact
and should not be reported as a fact about confidence.} That averaged
confidence is diluted by high-confidence filler is established rather than new:
it is the motivating premise of relevance-weighted uncertainty
estimation~\citep{duan2024shifting} and the target of recent length-invariant
estimators~\citep{fadeeva2025uncertaintyline}. The related pathology in machine
translation is that sentence-level model probability produces both a beam
problem and a brevity problem, which~\citet{murray2018correcting} trace to
label bias and correct with a sentence-level term, finding a per-word reward
slightly better than length normalization; that is a different mechanism from
ours and we cite it as a precedent for correcting a length-coupled score, not
as the same finding. Measuring at the answer span instead is likewise
established. \textbf{We adopt the fix rather than proposing it}, and its role
here is narrow: it closes off ``you measured in the wrong place'' as an
explanation for the next result.

\textbf{Signal 3: the answer-token margin. Fails on saturation, with a
mechanism.} Measured at the answer token, where dilution cannot apply, the
signal is saturated. Section~\ref{sec:commitment} gives the registered test.
What it supports, stated as an inference rather than an observation:
\emph{the answer token behaves as a near-deterministic readout of a chain that
has already committed}. We observe saturation downstream, at the answer
position, and infer commitment upstream; we do not observe the commitment
itself. The saturation is \emph{consistent with} a chain that fixed its answer
before the answer token was emitted, and it \emph{supports} that reading over
one in which the answer token is where the model decides, but it does not
establish when or where the decision was made. A sample whose answer
contradicts the plurality of its own problem still emits that answer at a
median 20.52 nats, odds of roughly 800 million to one against the nearest
alternative option. Dissenting samples are as committed as agreeing ones. On this
reading the variance self-consistency exploits is not expressed at the answer
token, and is more plausibly carried by which chain got written, though that
locus is inferred here and not measured. This is a mechanism rather than
a restatement because it is falsifiable and was registered in advance: had
dissenting samples been measurably less certain than agreeing ones, the
thresholds in Section~\ref{sec:commitment} would have failed.

\textbf{What this does to v1's thesis.} One of the three failures is not about
confidence, one is about confidence in a way v1 did not state, and one is
unexplained. \textbf{v1's ``confidence does not track correctness'' is retired
as a unifying thesis}, and this subsection must not be read as though every
failure now has an account. Two of three do.

\begin{figure}[t]
\begin{center}
\includegraphics[width=0.6\linewidth]{calibration_both.png}
\end{center}
\caption{Fraction of plurality answers correct by agreement bin. Both models
fall far below the calibration diagonal; Llama's highest-agreement bin is less
accurate than its lowest, with no monotone trend across bins.}
\label{fig:calibration}
\end{figure}

\begin{table}[t]
\begin{center}
\begin{tabular}{lcccc}
\toprule
\textbf{Confidence bin} & \textbf{Qwen $n$} & \textbf{Qwen frac correct} & \textbf{Llama $n$} & \textbf{Llama frac correct} \\
\midrule
{[}0.25, 0.50) & 56 & 33.9\% & 79 & 30.4\% \\
{[}0.50, 0.75) & 83 & 27.7\% & 84 & 33.3\% \\
{[}0.75, 1.00] & 59 & 52.5\% & 35 & 28.6\% \\
\bottomrule
\end{tabular}
\end{center}
\caption{Calibration by agreement bin (pooled 198). Confidence is the
plurality fraction over all samples.}
\label{tab:calibration}
\end{table}

\subsection{The commitment result, registered and held out}
\label{sec:commitment}

The margin claims were registered with thresholds fixed, then tested once on
the 69 problems no exploratory analysis had read.

\begin{table}[t]
\begin{center}
\begin{tabular}{llccc}
\toprule
\textbf{Id} & \textbf{Quantity, over dissenting samples} & \textbf{Estimate} & \textbf{95\% lower} & \textbf{Threshold} \\
\midrule
MD1 & per-problem median margin (nats) & 20.5232 & 18.8376 & 15.0 \\
MD2 & per-problem fraction above 10 nats & 0.7567 & 0.7105 & 0.60 \\
\bottomrule
\end{tabular}
\end{center}
\caption{Pre-registered descriptive claims on the 69 unexposed problems. Both
PASS. One-sided 95\% lower bounds, cluster bootstrap over problems. 64 of 69
problems carried at least three dissenting samples; 5 were excluded by the
registered rule and are counted, not imputed. Answer rate on the store is
0.9950 $[0.9936, 0.9961]$.}
\label{tab:md}
\end{table}

The holdout reproduces the exposed set closely, 20.52 against 20.62 and 0.757
against 0.749; only the holdout figures were registered.

\textbf{The two units, measured.} The margin is not devoid of information about
correctness, and the difference between the two units is not an argument we
make about the result: it is the result. Pooled at sample level over all 198
problems (4,283 correct and 8,322 incorrect samples with measured margins), the
fraction above 10 nats is 0.8228 for correct samples against 0.7624 for
incorrect, a difference of $+0.0604$ $[+0.0183, +0.1017]$, which excludes zero.
Per-problem and paired, over the 186 problems carrying at least one of each,
the mean difference is $-0.0168$ $[-0.0527, +0.0182]$, which crosses zero
(Figure~\ref{fig:twounits}). The population differs because a paired difference is undefined where a problem has
no correct or no incorrect sample: 2 problems are all-correct and 10
all-incorrect, and those 12 contribute to the pooled figure and not the paired
one.

\begin{figure}[t]
\begin{center}
\includegraphics[width=0.92\linewidth]{two_units.png}
\end{center}
\caption{The same margin, the same statistic, two units. \textbf{Pooled at
sample level} across all 198 problems (4,283 correct and 8,322 incorrect
samples carrying a measured margin), correct samples exceed 10 nats more often
than incorrect ones and the interval excludes zero: this answers ``does the
margin carry information about correctness across the benchmark?'', and the
answer is yes. \textbf{Per-problem and paired}, over the 186 problems carrying
at least one sample of each kind, the difference cannot be distinguished from
zero: this answers ``within a single problem, does the margin say which of
several disagreeing samples to trust?'', and the answer is no. The 12 problems
absent from the paired row are all-correct (2) or all-incorrect (10), where a
paired difference is undefined. \textbf{The two are not in conflict}: they are
different questions, and only the second is the one a self-consistency router
has to answer. Intervals are 95\% cluster bootstrap over problems.}
\label{fig:twounits}
\end{figure}

\textbf{We therefore do not claim that token log-probabilities are
uninformative}, and any reading of this paper reaching that conclusion has
overshot. Across questions the margin carries a real, small signal; within a
question it does not carry one we can detect, and within a question is where a
router stands. This makes the reconciliation with recent work a measurement
rather than an argument. \citet{kumaran2026confidence} reports that calibrated
log-probability confidence behaves as an answer-evidence signal coupled to
correctness, at AUROC 0.62 to 0.80. That quantity, a temperature-scaled softmax
over the option letters, is ours at the same locus. Its unit is not: every
result there is trial-level across questions, with one answer drawn per
question, and the design never conditions on samples that disagree with each
other because it never has two samples of one question to compare. Our pooled
figure reproduces the direction that paper reports; our paired figure is the
one it never measured.

\textbf{No margin gate was built, and this is not an evaluation of one.} We
measured the signal's distribution and report that it does not separate the
samples a within-question router would have to separate. \textbf{That is
evidence against the signal supporting the required decision, not an empirical
evaluation of a deployed margin gate}, and Section~\ref{sec:gates} accordingly
remains a two-gate result.

The objection this invites should be stated rather than left for a reader to
find: \textbf{a signal with weak paired discrimination can still support useful
routing if the routing rule is nonlinear} in it. A threshold that fires only in
the extreme tail, a rule combining the margin with agreement or with length, or
a per-problem rather than per-sample criterion could each extract value that a
paired mean difference does not see. None of those is tested here. What we
show is that the simplest reading, that a more certain sample is more likely to
be the right one within its own problem, does not hold; what we do not show is
that no rule built on this signal can work.

\section{Discussion}
\label{sec:discussion}

\textbf{Why backfire occurs.} On hard problems the model's sampling
distribution concentrates on a wrong answer, so voting locks in the error.
Backfire is not a sampling artifact; it reflects the per-problem answer
distribution. GPQA distractors are designed to be
plausible~\citep{rein2023gpqa}, which amplifies the effect.

\textbf{Why the gates fail (revised in v2).} v1 wrote: ``Both gates read
confidence (agreement, or low entropy), but confidence on these problems does
not indicate correctness.'' That sentence is withdrawn as stated. The entropy
gate was not reading confidence; it was reading a diluted average dominated by
prose. The margin, which does read confidence at the right place, fails because
the answer token is saturated whether or not the sample agrees with its own
plurality. And the agreement gate's failure is explained by neither account.
Three signals, three statuses, set out in Section~\ref{sec:why}.

\textbf{Positioning.} Prior work established that self-consistency helps on
benchmarks where models have higher baseline
accuracy~\citep{wang2023self}; we show it hurts the majority of problems on a
hard one. Repeated-sampling analyses study coverage under an oracle
metric~\citep{brown2024large}; we instead decompose the gap between realizable
majority vote and the oracle. Adaptive-consistency early
stopping~\citep{aggarwal2023adaptive} is essentially our agreement gate,
validated on easier datasets where backfire is rare; our negative result shows
agreement stability is insufficient on hard problems. Closest to our work,
\citet{chen2024more} show that majority-vote accuracy can rise and then fall
as the number of LM calls grows, attribute this non-monotonicity to a mixture
of easy and hard queries within a task (more calls help the easy ones and hurt
the hard ones), and use that structure to estimate, from a small number of
samples, the call count that maximizes aggregate performance. We differ in two
ways. We measure the fraction of individual problems harmed rather than the
shape of the aggregate curve, which matters because aggregate accuracy here
looks nearly flat (0.342 to 0.369 for Qwen) while a majority of problems
degrade underneath it. And our gate results sit in tension with few-sample
estimation of an optimal call count. v1 put this as a claim about a uniformly
hard benchmark; the stored data does not support that description for these
models (disclosure 10), so the claim is narrower than v1 stated. On this
benchmark, whose per-problem difficulty is measurably mixed rather than
uniform, the
verifier-free signals such estimation would rely on do not recover the
headroom. \citet{tan2025consistent} study self-consistent errors, where a
model repeats the same wrong answer across samples, and find that
consistency-based detectors, of which semantic cross-checking is one
example~\citep{zhang2023sac3}, fail on exactly those cases, which leads them to
argue for an external verifier; their setting is error detection rather than
compute allocation, but the conclusion converges with ours. The external-signal
direction this motivates, trained verifiers and process reward
models~\citep{cobbe2021training, lightman2023verify}, is the natural way to
reach the oracle headroom, which our two verifier-free gates cannot.

\textbf{Implications.} On hard inputs, do not assume self-consistency is safe,
and do not expect agreement or token entropy to tell you when it is; we tested
both and both fail. Recovering the headroom likely requires a signal external
to the model's own samples.

\section{Limitations}

\textbf{Llama is near chance.} Llama-3-8B's single-sample accuracy on full
GPQA Diamond is 0.273, just above the 0.25 random baseline, so its backfire
is less surprising. Qwen (0.342, clearly above chance) is the stronger
demonstration; Llama corroborates the direction.

\textbf{One benchmark.} All results are on GPQA Diamond, graduate-level
science. Other domains and easier difficulty regimes may differ.

\textbf{Two small, non-reasoning models (rescoped in v2).} Reasoning-native models require
different infrastructure. We attempted a preliminary evaluation of a
reasoning-native model with hidden chain-of-thought and found that a
substantial fraction of samples exhausted the output token budget on hidden
reasoning before producing a visible answer. This truncation made
majority-vote accuracy incomparable to our non-reasoning models. A proper
evaluation of reasoning-native architectures therefore requires larger
inference budgets and explicit handling of hidden CoT output length, and is
left to future work. Whether backfire shrinks or persists for such models
remains the central open question. \textbf{This limitation is no longer open as
stated.} Section~\ref{sec:wall} reports the attempt: three models and three
separately measured walls. On hosted serverless inference at a small budget the
question is not measurable, and it remains open for the open-weights route.

\textbf{Subset representativeness.} The original 47-problem subset was easier
than the full benchmark (Qwen $N=1$ accuracy 0.418 vs 0.342 pooled), which is
why we report on all 198.

\textbf{Monte Carlo boundary sensitivity.} $\text{MV\_acc}(N)$ is estimated by
Monte Carlo, so a problem whose $\text{mv\_gain}$ is exactly zero can fall on
either side of the backfire boundary between runs. The pre-registration records
Qwen's exploratory backfire rate as 46.8\% (22 of 47); recomputing it in the
pipeline that produced the results reported here gives 44.7\% (21 of 47), a
one-problem difference. Llama's exploratory rate is unchanged at 66.0\%. The
pre-registered prediction is unaffected, since both values sit far above the
33\% threshold, but the same sensitivity applies in principle to the
confirmatory rates: 14 of Qwen's 151 confirmatory problems have exactly zero
gain. Because backfire is defined strictly as $\text{mv\_gain} < 0$, those 14
currently count as non-backfire, so the sensitivity is one-sided and the rate
can only rise. Even if all 14 crossed the boundary, Qwen's confirmatory rate
would be 69.5\% rather than 60.3\%, still far above the 33\% threshold. We
report the recomputed values throughout and leave the pre-registration as
tagged.

\textbf{Oracle and ratio.} The oracle is an upper bound, not deployable, and
the fraction-of-oracle-captured ratio is noisy (small denominator), so we lead
gate claims with absolute accuracy.

\textbf{Entropy threshold selection.} The entropy threshold was chosen by an
in-sample argmax over 40 candidate values on the confirmatory 151, maximizing
capture against the binary $\{N=1, N=64\}$ oracle relative to an
$\text{MV\_acc}(64)$ baseline, whereas the capture figures we report in
Section~\ref{sec:gates} are measured against the grid oracle relative to
$\text{MV\_acc}(1)$. Because selection and evaluation use the same problems,
the entropy gate's reported capture is optimistic.

\textbf{Limited gate family (resolved in v2).} v1 evaluated two verifier-free
signals and noted that richer signals such as final-answer log-probability
margins could behave differently, but that evaluating them requires per-token
log-probability arrays which its pipeline did not store. This version stores
them and evaluates that signal in Sections~\ref{sec:why}
and~\ref{sec:commitment}. It does not behave differently in the way that
sentence hoped.

\textbf{One model for the new claims.} The registered second-model replication
was sampled and could not be evaluated (Section~\ref{sec:rejected}).
Cross-family replication was never purchasable either: four of five priced
candidates in the 7 to 9B range refused serverless requests.

\textbf{Agreement's failure has no mechanism.} See Section~\ref{sec:why}. This
is the largest unresolved question in the paper, and v2 makes it more visible
rather than smaller.

\textbf{The reasoning wall is a serverless result.} All three models in
Section~\ref{sec:wall} are downloadable, and on fixed-cost compute the caps
that defeated us are affordable. That section bounds what a metered API buys at
this budget, not what is knowable.

\textbf{No compute-matched baseline}, in v1 or v2. Every gate comparison is
against a flat $N=64$.

\textbf{No margin gate was built}, so the margin's failure is inferred from the
signal's distribution rather than measured as a routing outcome.

\textbf{An exploratory length--accuracy shape is not established.} Binning
problems by per-problem mean completion length gives accuracies of 0.4441,
0.3400, 0.2218, 0.3316 and 0.4098 across quintiles, but the shape-agnostic
quadratic term gives $p = 0.083$. Two mechanisms were ruled out, truncation
selection and within-problem bimodality. It was found on an already-exposed set
with its sharpest contrast chosen after inspection, and it appears here rather
than in results for that reason.

\textbf{Known mechanism, narrowed in v2.} Confidence failing to track
correctness is an instance of documented
overconfidence~\citep{guo2017calibration, kadavath2022language}. v2 narrows
what we claim from it: only one of the three signal failures is explained that
way, and the unit matters (Section~\ref{sec:commitment}). Our contribution is
the pre-registered, replicated, self-consistency-specific quantification and
the failure of three cheap signals, not the existence of miscalibration.

\section{The reasoning wall, on hosted serverless inference}
\label{sec:wall}

v1's Limitations reported that a preliminary reasoning-native evaluation had
samples exhausting the output budget on hidden reasoning, and left the question
open. This section reports what happened when we tried to close it.

\textbf{Scope, stated before the measurements.} The claim here is about hosted
serverless inference and does not hold for open weights on owned or free
compute. All three models below are downloadable. Nothing here says a
reasoning-native model cannot be evaluated on GPQA Diamond; it says that on a
metered per-token hosted API at a small budget, three of them could not be.
Serverless pricing charges per output token and reasoning-native models emit a
great many, so the cost of a fixed experiment scales with the model's verbosity
rather than with the size of the benchmark. The wall is a property of the
billing model meeting the generation length, not of the models being hard to
run. The route a follow-up should take is open weights on fixed-cost compute,
where the marginal cost of an output token is zero and the binding constraint
becomes wall-clock and memory: a free tier offering two 16\,GB GPUs and roughly
30 GPU-hours per week is enough for a 7 to 9B model in reduced precision at
long generations, though not for a 32B model without quantization or offload.
Under that budget the caps that defeated us are affordable. What changes is not
the science but the denominator.

\textbf{Position relative to prior work.} That token budgets change evaluation
outcomes is established. \citet{souza2026budget} vary the generation budget
across seven levels from 64 to 4,096 tokens, over four models and three
benchmarks including GPQA Diamond at the same 198 items, 56,476 inferences in
total, and report that model rankings reverse across budgets on all benchmarks
($p < 0.01$, McNemar), with 3 to 19\% of items non-monotone even after
controlling for truncation through a three-tier analysis. That work is stronger
than this on the general claim. Two things about it define this section's
scope. First, it deliberately excludes \emph{dedicated} reasoning models,
naming o1, DeepSeek-R1 and QwQ, while describing its own four as open-weight
reasoning models: the exclusion is specifically of hidden chain-of-thought
architectures, which is exactly our three. Second, its budget grid stops at
4,096 tokens, and one of our models returns an answer rate of 0.0000 at both
2,048 and 4,096, so across that entire grid the measurement this section
reports does not exist at all.

\textbf{Three models, three measured walls}, each established by a real request
rather than inferred from a price list. QwQ-32B is unreachable:
\texttt{model\_not\_available}, no serverless route at any price. MiniMax-M2.7
answers, but never at a comparable cap: answer rate 0.6460 at cap 16,384 and
0.2649 at the 2,048 v1 used. Qwen3.5-9B returns nothing at any affordable cap:
answer rate 0.0000 at both 2,048 and 4,096, with \emph{mean completion equal to
the cap exactly}, and 0.5938 at 8,192. The reference is
Qwen2.5-7B-Instruct-Turbo at cap 2048, which answers 0.9950 over 12,672
samples. Every one of 52 samples at 2,048 and 4,096 ran to the ceiling, and at
4,096 the visible channel received 38 characters. That model's own card
recommends 32,768 output tokens for general queries, so v1's 2,048 is a
sixteenth of this model class's lower recommendation: this is not a model
failing at a reasonable cap, it is a cap fixed before this class of model was
the default.

\textbf{The door, and what it changes.} The provider honours a thinking control
on Qwen3.5-9B. Two spellings are accepted and indistinguishable in effect: mean
completion 1537.3 against 1532.8 at cap 2,048, identical truncation, no
reasoning field returned either way. But a reasoning model with reasoning
disabled is a second non-reasoning model. What became purchasable was not the
replication that was wanted, and we do not describe it as one.

\textbf{Comparability is keyed on answer rate, not on matched caps.} Two models
at the same cap with different length distributions produce two different
output populations wearing one benchmark's name. MiniMax-M2.7 at 2,048 answers
0.2649; comparing its answered samples against Qwen's near-complete ones
compares two different sets of problems. Two conditions are comparable when
their answer rates match, and the answer rate is published beside every
accuracy. This differs from post-hoc truncation filtering in when it applies:
it is a design constraint on what to sample, not a repair applied to an
existing comparison.

\textbf{Cost and accounting.} What this cost, how the ledger reconciles
against the provider's own billing export, and how a small probe's cost
projection should be calibrated are reported in
Appendix~\ref{app:accounting}. The scientific claim of this section does not
depend on them: it is that under the serverless budget available to us, these
models could not be evaluated on a footing comparable to the non-reasoning
models, for the three separately measured reasons above.

\subsection{A gated benchmark cannot have its chains published}
\label{sec:cot-release}

An infrastructure constraint we did not anticipate and found by running the
check rather than by reasoning about it. Policy in this project is that
benchmark question text is never committed: the sample record stores a prompt
hash and never the prompt. A pre-release leakage check reduces the gated
question source to hashed $n$-grams and scans the release tree for them. Run
over the raw store it reports 5,669 hits across 198 files and blocks the
release; run over the full raw release tree it reports 20,376. The cause is not
a stored prompt. It is the completion text: \textbf{the model restates the
question inside its own chain of thought.} The request never carried the
question into storage; the response did. The consequence generalizes past this
project. A chain-of-thought sample store on a gated benchmark cannot be
published verbatim, however careful the storing pipeline is, because the leak
is authored by the model rather than by the pipeline, and redaction rules
written for what a project writes do not cover what the model writes back. This
sits awkwardly beside reproducibility: the raw store is what regenerates every
derived number, and it is the artifact that cannot be released. Our resolution
is the two-part release recorded in the Data availability statement, a
derived-only open artifact and a gated raw one on the same access terms as the
benchmark. We flag it because every paper that publishes reasoning traces on a
gated benchmark faces it, and we have not seen it stated.

\section{When the comparability rule rejected our own registered result}
\label{sec:rejected}

The second-model replication was registered and is \textbf{not evaluated}. Two
claims restated the Section~\ref{sec:why} mechanism keyed on correctness rather
than plurality, with thresholds 15.0 nats and 0.60 set below the exposed
estimates of 21.5365 and 0.7617 so that each was a prediction, and the
registration was tagged before any confirmatory sample existed. A stratification
check found no dependence of either quantity on per-problem accuracy, so the
thresholds were held rather than recalibrated.

\textbf{Cap selection failed, and not randomly.} A 32-sample probe at cap 6,144
measured answer rate 1.0000 and truncation 0.0000. The run measured 0.8931
$[0.8747, 0.9090]$ and truncation 0.2905, with mean completion 65.3\% above
projection. Configuration drift was ruled out before the selection explanation
was accepted: one parameter hash across all 1,253 samples matching the
thinking-disabled configuration, zero nonzero reasoning-token counts, cap 6,144
throughout. The cause is measurable because the sampler iterates problem-major
and so re-sampled the probe's own problems first: those 8 problems give mean
completion 2082.9 and truncation 0.0312 across 128 samples, while every other
problem reached gives 3546.4 and 0.3200 across 1,125. The probe reproduced
itself and was precise about an unrepresentative slice, since the rest of the
benchmark runs 1.703 times longer. The $k=8$ multiplier of 22.7\% could not
cover a 65.3\% error, because that figure describes a \emph{random} draw of 8
problems and a fixed lowest-id slice is not one: its error is not resampleable,
because there is one such slice and it is the same every time.

\textbf{Consequence.} At 29\% truncation the pool that would be scored is
selected for finishing fastest, and those samples are missing non-randomly.
Scoring them would produce a number about the subset of samples that fit inside
6,144 tokens and report it as a number about the model, which is the confound
the comparability rule exists to prevent. The two claims are therefore
registered and unevaluated: not withdrawn, not falsified. Nor was the
replication recoverable at a different cap. A right-censored fit to the 1,253
samples puts the cap required for the registered 0.9950 answer rate near 41,000
output tokens and the corresponding run at about \$4.19, against \$1.16
remaining; even a 0.95 rate, which would still fail the condition, requires
roughly 16,268 tokens and \$3.84. The replication was not lost by choosing
6,144. It was not purchasable at any cap on this budget.

\textbf{The point.} A methodological rule that never rejects anything is
decoration. This one rejected a result its own authors had registered, sampled
and paid for.

\section{A reconstructed few-sample estimator}
\label{sec:threada}

\textbf{The premise was tested first, and the first formulation failed it.}
An estimator that infers an optimal call count from a few samples relies on the
task containing a mixture of easy and hard queries: without that mixture there
is nothing for it to detect. We therefore tested whether GPQA Diamond carries
one for these models before testing the estimator on it. It does.
\textbf{Per-problem single-sample accuracy across all 198 problems varies at
25.5 times a homogeneous null for Qwen and 17.1 times for Llama}, $p < 0.0001$
in each case over 10,000 permutations, with 45 and 37 problems respectively
sitting at or below 0.05 or at or above 0.95 against a null expectation near
zero. GPQA Diamond is hard on average and is \emph{not} uniformly hard for
these models: it carries exactly the easy/hard mixture~\citet{chen2024more}
rely on. The estimator is therefore tested where the mixture is present and
measured, rather than where its premise is absent, and a negative result here
is informative about the estimator rather than about the benchmark.

\textbf{Mixture presence turned out to be necessary but not sufficient.} The
mixture is present, and the non-monotonicity it is supposed to produce is not:
the confirmatory aggregate accuracy curve across the grid spans 0.52 points and
does not fall. Whatever makes an aggregate curve rise and then decline, the
presence of a per-problem difficulty mixture does not guarantee it, at least at
this benchmark's difficulty and this scale of model.

We reconstructed the few-sample call-count estimator that follows
from~\citet{chen2024more} and pre-registered a comparison against a naive
within-$k$ baseline on v1's confirmatory 151, using paired per-problem regret
where negative favours the estimator. At $k=8$ the difference is $-0.0213$
$[-0.0335, -0.0091]$, a PASS by more than the registered resolution floor of
0.0161. At $k=4$ it is $-0.0186$ $[-0.0332, -0.0041]$, a PASS that was
registered as underpowered. At $k=16$ it is $-0.0040$ $[-0.0085, +0.0004]$, a
\textbf{FAIL} by 0.0004, reported as a failure rather than rounded into the
pattern of the other two. A shrinkage baseline was added post hoc and is
labelled as such: at $k=4$ and $k=16$ every shrinkage strength is worse than
the plain baseline, and at $k=8$ the weakest helps slightly. It is not
registered and decides nothing.

\section{Corrections and disclosures}
\label{sec:corrections}

Recorded as a section rather than a footnote.

\begin{enumerate}
\item \textbf{A superseded draft was cited} in place of the published preprint
in project documentation. A test now scans every document in the repository for
citations to superseded sources.
\item \textbf{Three per-problem properties were described as independent}
before being tested. Within one model at $n=127$, per-problem accuracy against
the sub-2-nat margin tail gives $-0.2614$ $[-0.3912, -0.1156]$, excluding zero.
\item \textbf{The margin run changed concurrency mid-run.} Because the sampler
iterates problem-major, the resulting regime comparison is between problems
rather than within them, so it is confounded and decides nothing. Three
manifests are marked as reconstructed rather than contemporaneous.
\item \textbf{An option-order claim was corrected after tagging.} Options are
shuffled by a row-index-keyed generator, reproducing v1's shuffle; prompt
hashes are equal across both model stores for all 198 problems, verified.
\item \textbf{An answer-rate pairing rule did not cover prose tables} until a
quintile table of accuracies was published without them.
\item \textbf{The second-model cap was chosen from 8 non-random problems}
(Section~\ref{sec:rejected}). Probe problems are now drawn at random with a
recorded seed.
\item \textbf{v1's own split has a leak:} 3 of its 50 exploratory ids sit
inside the confirmatory 151. Disclosed here; v1's verdicts are not recomputed.
\item \textbf{A bibliography entry's authors were written from memory.} The
entry for a length-invariant uncertainty estimator was added from a
search-result title rather than from the paper: its author list named a person
who is not an author and omitted two who are. It survived a merge, a build and
a full pre-submission read, and was caught in an entry-by-entry audit asking
which references had actually been opened. Corrected against the arXiv record.
\textbf{The root cause is the one this paper has disclosed twice already}: a
claim made from recollection rather than from a source, of the same kind as
citing a superseded draft and as asserting three properties were independent
before testing them. It is now a check rather than a habit: every entry
carrying an arXiv id is compared against the arXiv API for title, full author
list and year, and every entry without one must record which source was opened
and when, with an entry lacking that record failing the suite. All sixteen
entries, including the eleven inherited from v1, were re-verified this way; no
further mismatch was found.

\item \textbf{v1's Setup describes uniform-random tie-breaking; its code breaks
ties lexicographically.} Found while reproducing v1's figures from the archived
derived data: the plurality function sorts the tied answers and returns the
first. The published numbers follow the code. Quantified rather than only
stated: Qwen has 5 full-pool plurality ties of 198 and Llama 2, and counting
$N=64$ subsets Qwen has 47 where the correct answer is in a tie while
\textbf{Llama has none}, so Llama's figures are convention-independent
outright. Lexicographic order is not a biased choice on this dataset: the
correct answer is A on 48 problems, B on 51, C on 47 and D on 52 after the
shuffle, $\chi^2 = 0.343$ on 3 degrees of freedom. Across 5,000 uniform-random
tie-break seeds Qwen's $\text{MV\_acc}(64)$ has mean 0.3726 and standard
deviation 0.0051, spanning 0.3619 to 0.3833, and the lexicographic 0.3686 sits
at the 31.7th percentile: inside the distribution, slightly unlucky, not an
outlier. Backfire spans 109 to 113 of 198. \textbf{The convention is a
third-decimal effect and changes no conclusion:} every seed leaves backfire far
above the 33\% pre-registered threshold and still a majority, and every seed
leaves $\text{MV\_acc}(64)$ above $\text{MV\_acc}(1)$ of 0.3419. v1's verdicts
are not recomputed.

\item \textbf{v1's Positioning describes the benchmark as uniformly hard; the
stored data does not support that for these models.} Per-problem single-sample
accuracy varies at 25.5 times a homogeneous null for Qwen and 17.1 times for
Llama, $p < 0.0001$ over 10,000 permutations, with 45 and 37 problems at or
below 0.05 or at or above 0.95. GPQA Diamond is hard \emph{on average} and is
not uniform in difficulty for these models. \textbf{This touches no v1 result
and no v1 number}: it bears only on the positioning against~\citet{chen2024more},
whose account turns on exactly that mixture. It is disclosed rather than quietly
corrected because Section~\ref{sec:threada} builds on the same positioning, and
using a premise in one section while leaving its overstatement standing in
another is the kind of quiet reliance this section exists to prevent. The
sentence in Section~\ref{sec:discussion} is qualified accordingly.
\end{enumerate}

\section{Conclusion}

On hard reasoning problems, self-consistency backfires on the majority of
problems, pre-registered and confirmed on the full GPQA Diamond benchmark for
Qwen and corroborated by a second family from a near-chance baseline.
\textbf{That result is unchanged from v1.}

What has changed is the explanation for why cheap verifier-free gates cannot
recover the headroom. v1 attributed it to confidence not tracking correctness.
With a third signal measured at the right place, that account no longer covers
the evidence. The entropy gate was reading a diluted statistic rather than a
confidence signal. The answer-token margin, read where dilution cannot apply,
fails because the answer token is saturated: a sample that contradicts its own
problem's plurality is as committed as one that agrees with it, so the variance
a router would need has already been spent upstream. And the agreement gate's
failure is explained by neither, and remains open. Recovering the headroom
likely requires a signal external to the model's own samples, or one read
earlier in the chain than the answer. Whether reasoning-native models escape
any of this remains open, and on hosted serverless inference at a small budget
it is not measurable at all.

\appendix

\section{Cost, accounting, and probe calibration}
\label{app:accounting}

Lifted out of Section~\ref{sec:wall} so that the body carries the scientific
claim and this carries the bookkeeping. Nothing is omitted.

\textbf{The bill, reconciled against the provider.} Our client-side ledger
records \$4.697153 across 14,125 sampled requests. The provider's own billing
export, filtered to this project's API key and summed from unrounded quantities
times unit prices rather than the rounded per-line amounts, gives
\textbf{\$4.701169} across six line items over two days. The difference is
\textbf{\$0.004017, or 0.085\%}, and the ledger is the lower of the two.
\citet{souza2026budget} report 56,476 API calls and no monetary cost. We report
both figures because the negative results in this paper are only interpretable
alongside what was affordable, and because a cost claim that is not reconciled
against the counterparty is an assertion rather than an audit.

\textbf{Where the 0.085\% goes, and why it is one-sided.} The residual is
about 4,600 unrecorded tokens on one model and 11,000 on the other. A
client-side ledger writes a row when a response \emph{arrives}, so it cannot
see a request that was issued, billed, and then abandoned. This run has three
such boundaries: a mid-run concurrency change, a \texttt{SIGTERM} that stopped
the second-model run at concurrency 16, and an interrupted probe. Each leaves
its in-flight requests billed by the provider and absent from the ledger.
\textbf{This is a property of the accounting method rather than a discrepancy
to explain away}: a client-side ledger undercounts any interrupted run by
exactly its in-flight requests, always in the same direction, and the size of
the undercount is the concurrency times the number of interruptions. Anyone
reporting costs this way should expect a small one-sided gap and should say so
rather than round it off.

\textbf{The reconciliation is only possible because the key was dedicated.}
Project security policy required a separate API key per project, for revocation
rather than for accounting. The same billing account carries a second key with
\$0.292479 of unrelated spend on 10 and 11 August, including a model this
project never used, and all twelve line items across both keys sum to \$4.99.
With a shared key the Argmax bill would not have been separable from that at
all: no filter, no reconciliation, and the \$4.99 would have had to be reported
as the project's cost or not reported. \textbf{A security rule turned out to be
an accounting rule}, and we would not have discovered that had the second key's
spend not happened to overlap in time.

\textbf{Projecting cost from a small probe.} Mean completion length is a
per-problem property at 119.66 times the variance a homogeneous null produces,
so a probe over $k$ problems inherits that between-problem spread and taking
more samples per problem does not reduce it. Resampling the 198 per-problem
means, 20,000 trials per $k$, the multiplier needed for a projection to cover
the truth 95\% of the time is 34.9\% at $k=4$, 22.7\% at $k=8$, 15.2\% at
$k=16$ and 10.0\% at $k=32$. Probes are near-unbiased in the median, yet just
over half underestimate at every $k$ because per-problem means are
right-skewed, so a ceiling set at the projection is wrong about half the time.


\bibliography{references}
\bibliographystyle{colm2026_conference}

\end{document}
